\documentclass[11pt,a4paper]{article}  % Déclare la classe du document.

\newcommand{\chaptermark}[2]{\og #2 \fg{} (#1)} % ligne uniquement pour éviter une erreur en cas d'article
\include{../1ere-preambule}

\usepackage{epic}   % extension et simplification de  l'environnement picture
\usepackage{graphicx}    % permet d'inclure des graphique externe dans un document
\usepackage{arydshln}   %permet d'avoir des lignes en pointillés

\newcommand{\lignes}[1]{
	\setlength{\unitlength}{1mm}
	\vskip 0,4cm\noindent
	\multido{\i=0+1}{#1}{%
		\noindent
		\begin{picture}(150,5)
			\linethickness{0,005mm}
			\put(-5,0){\dashline{1}(175,5)(5,5)}
		\end{picture}\vskip 0,001cm
	}
	\vskip 0cm
}

\rhead{Chapitre 01 - Fonctions de référence - Exercices d'entrainements}

\newcommand{\va}[1]{\lvert \ #1 \ \rvert}
\newcommand{\vaf}[1]{\left\lvert \ #1\  \right\rvert}

\begin{document}
	\pagestyle{empty}
	\graphicspath{{images/}}
	
	\vspace*{-2cm}
	
	\begin{center}
		\fbox{
			\begin{minipage}[]{15cm}
				\begin{center}
					\LARGE{\rule{0.5cm}{0pt}\rule[-0.25cm]{0pt}{0.9cm}
						Déterminer une équation de droite
					}
				\end{center}
			\end{minipage}
		}
	\end{center}

\begin{methodeTitre}{Déterminer une équation réduite de droite dont on connaît deux points}
	\grasInMet{Énoncé :}\\
	Soit \Rij un repère du plan. Soit A\coord{4}{-1} et B\coord{3}{5} deux points d'une droite $d$.
	Déterminer une équation de la droite $d$.\\
	\ \\
	\grasInMet{Réponse :}\\
	Les points A et B sont d'abscisses différentes donc la droite $d$ n'est pas parallèle à l'axe des ordonnées.\\ Elle est donc de la forme $y = mx + p$, où $m$ et $p$ sont deux nombres réels.\\
	Le coefficient directeur de $d$ est $m =  \dfrac{y_B-y_A}{x_B-x_A}=\dfrac{5-(-1)}{3-4}=\dfrac{6}{-1} = -6$.\\
	L'équation de $d$ est donc de la forme : $y = -6x + p$\\
	\grasInMet{Soit :}\\
	Comme A\coord{4}{-1} appartient à la droite $d$, ses coordonnées vérifient l'équation de $d$ soit :\\
	$-1 = -6 \times 4 + p$. D'où  $p = -1 + 6 \times 4 = 23$\\
	\grasInMet{Soit :}\\
	Comme B\coord{3}{5} appartient à la droite $d$, ses coordonnées vérifient l'équation de $d$ soit :\\
	$5 = -6 \times 3 + p$. D'où  $p = 5 + 6 \times 3 = 23$\\
	\ \\
	Une équation de $d$ est donc :  $y = - 6x + 23$.
\end{methodeTitre}
\ \\
\ \\
\ \\
	\begin{center}
		\fbox{
	\begin{minipage}[]{15cm}
		\begin{center}
			\LARGE{\rule{0.5cm}{0pt}\rule[-0.25cm]{0pt}{0.9cm}
				Déterminer une équation de droite
			}
		\end{center}
	\end{minipage}
}
\end{center}

\begin{methodeTitre}{Déterminer une équation réduite de droite dont on connaît deux points}
\grasInMet{Énoncé :}\\
Soit \Rij un repère du plan. Soit A\coord{4}{-1} et B\coord{3}{5} deux points d'une droite $d$.
Déterminer une équation de la droite $d$.\\
\ \\
\grasInMet{Réponse :}\\
Les points A et B sont d'abscisses différentes donc la droite $d$ n'est pas parallèle à l'axe des ordonnées.\\ Elle est donc de la forme $y = mx + p$, où $m$ et $p$ sont deux nombres réels.\\
Le coefficient directeur de $d$ est $m =  \dfrac{y_B-y_A}{x_B-x_A}=\dfrac{5-(-1)}{3-4}=\dfrac{6}{-1} = -6$.\\
L'équation de $d$ est donc de la forme : $y = -6x + p$\\
\grasInMet{Soit :}\\
Comme A\coord{4}{-1} appartient à la droite $d$, ses coordonnées vérifient l'équation de $d$ soit :\\
$-1 = -6 \times 4 + p$. D'où  $p = -1 + 6 \times 4 = 23$\\
\grasInMet{Soit :}\\
Comme B\coord{3}{5} appartient à la droite $d$, ses coordonnées vérifient l'équation de $d$ soit :\\
$5 = -6 \times 3 + p$. D'où  $p = 5 + 6 \times 3 = 23$\\
\ \\
Une équation de $d$ est donc :  $y = - 6x + 23$.
\end{methodeTitre}

\end{document}

